\chapter{Valorisation et états financiers}\label{ch:iv2:valuation-and-financial-statements}

\begin{ivtrains}
L'entretien de l'analyste fondamental~: suivre une transaction à travers les trois états financiers, construire une valorisation rapide à la main, dire ce que suppose un multiple, lire les résultats pour ce qu'ils cachent, et défendre une position quand les chiffres bougent. Chaque énoncé fournit tous les chiffres nécessaires. Travailler sur papier, dire ce que signifie chaque étape et terminer par ce que l'on ferait.
\end{ivtrains}

\begin{question}[$\star$ \iqroles{Analyste} \iqfirm{hedge fund, banque} \iqlevel{junior}]\label{iq:iv2:valuation-and-financial-statements:1}
La dotation aux amortissements d'une entreprise augmente de 10 pour l'exercice, et le taux d'impôt est de 25~\%. Cette dotation est déductible pour le même montant au titre de l'impôt, et rien d'autre ne change. Détaillez le compte de résultat, le tableau des flux de trésorerie et le bilan. De combien varient le résultat net, le flux de trésorerie opérationnel, la trésorerie, les immobilisations et les bénéfices non distribués~?
\end{question}

\begin{question}[$\star$ \iqroles{Analyste} \iqfirm{hedge fund, banque} \iqlevel{junior}]\label{iq:iv2:valuation-and-financial-statements:2}
Une nouvelle entreprise achète des marchandises pour 60 à crédit et les revend 100 à crédit. En fin d'exercice, elle n'a payé ni son fournisseur ni encaissé auprès de son client. Elle doit un impôt de 25~\% sur son bénéfice et le règle en trésorerie au cours de l'exercice. Quels sont son résultat net et son flux de trésorerie opérationnel~? Montrez que le bilan est équilibré.
\end{question}

\begin{question}[$\star$ \iqroles{Analyste, Trader} \iqfirm{hedge fund, banque} \iqlevel{junior}]\label{iq:iv2:valuation-and-financial-statements:3}
Une entreprise a une capitalisation boursière de 1\,000, une dette de 300, une trésorerie de 100, des actions de préférence pour 50 et des intérêts minoritaires de 30, ainsi qu'un EBITDA, c'est-à-dire son résultat avant intérêts, impôts et dotations aux amortissements, de 160. Quelles sont sa valeur d'entreprise et son multiple VE/EBITDA~? La valeur d'entreprise change-t-elle si elle emprunte 100 et conserve la trésorerie, verse un dividende de 50, ou dépense 100 de trésorerie en rachat d'actions~?
\end{question}

\begin{question}[$\star$ \iqroles{Analyste} \iqfirm{hedge fund, banque} \iqlevel{junior}]\label{iq:iv2:valuation-and-financial-statements:4}
Une entreprise a un résultat opérationnel de 200, un taux d'impôt de 25~\%, des dotations aux amortissements de 40, des dépenses d'investissement de 70, et un besoin en fonds de roulement qui augmente de 20. Elle verse des intérêts de 20 et emprunte 10 de plus, en net. Quels sont son flux de trésorerie disponible pour l'entreprise et son flux de trésorerie disponible pour les actionnaires~? Lequel des deux actualisez-vous au coût moyen pondéré du capital~?
\end{question}

\begin{question}[$\star$ \iqroles{Analyste, QR} \iqfirm{hedge fund, banque} \iqlevel{junior}]\label{iq:iv2:valuation-and-financial-statements:5}
Une entreprise produira un flux de trésorerie disponible de 100 l'an prochain, en croissance de 3~\% par an à perpétuité, et son coût du capital est de 9~\%. Quelle est sa valeur~? Quelle serait-elle si la croissance était de 4~\%, ou si le coût du capital était de 8~\%~? Qu'est-ce qui change si les 100 correspondent au flux de cette année plutôt qu'à celui de l'an prochain~?
\end{question}

\begin{question}[$\star$ \iqroles{Analyste} \iqfirm{hedge fund, banque} \iqlevel{junior}]\label{iq:iv2:valuation-and-financial-statements:6}
Un titre se négocie à 25, avec 100 actions en circulation. Il existe aussi 10 options de prix d'exercice 10, 2 actions soumises à restriction qui n'ont pas encore été acquises, et une obligation convertible d'une valeur nominale de 200 qui se convertit à 20 par action. Combien y a-t-il d'actions après dilution, et quelle est la valeur des fonds propres dilués~? Que faites-vous des 200 de l'obligation convertible lorsque vous passez de la valeur des fonds propres à la valeur d'entreprise~?
\end{question}

\begin{question}[$\star$ \iqroles{Analyste, Trader} \iqfirm{hedge fund, banque} \iqlevel{junior}]\label{iq:iv2:valuation-and-financial-statements:7}
Une entreprise réalise un chiffre d'affaires de 100, des coûts variables égaux à 20~\% du chiffre d'affaires, des coûts fixes de 60 et des intérêts de 10. Si le chiffre d'affaires augmente de 10~\%, de quel pourcentage augmentent le résultat opérationnel et le résultat avant impôt~? Pour quel chiffre d'affaires le résultat opérationnel s'annule-t-il, et pour quel chiffre d'affaires le résultat avant impôt~?
\end{question}

\begin{question}[$\star\star$ \iqroles{Analyste, QR} \iqfirm{hedge fund, banque} \iqlevel{junior}]\label{iq:iv2:valuation-and-financial-statements:8}
Une entreprise a des fonds propres valant 800 et une dette valant 200. Son bêta des fonds propres est de 1.2, le taux sans risque est de 4~\%, la prime de risque actions est de 5~\%, son coût de la dette est de 5~\% et son taux d'impôt est de 25~\%. Quel est son coût moyen pondéré du capital~? Elle passe ensuite à des montants égaux de dette et de fonds propres en valeur, et son coût de la dette monte à 6~\%. En retirant puis en rétablissant l'effet de levier avec la formule qui suppose que la dette a un bêta nul, quels sont son nouveau bêta des fonds propres et son nouveau coût moyen pondéré du capital~?
\end{question}

\begin{question}[$\star\star$ \iqroles{Analyste, QR} \iqfirm{hedge fund, banque} \iqlevel{junior, confirmé}]\label{iq:iv2:valuation-and-financial-statements:9}
Une entreprise sans dette gagne 100 après impôt chaque année à perpétuité et la distribue intégralement. Son coût du capital est de 10~\%, et le taux d'impôt est de 25~\%. Elle emprunte 500 à perpétuité à 5~\%, taux sans risque, et utilise la somme pour racheter des actions. Supposez que les intérêts sont déductibles de l'impôt et qu'il n'existe aucun coût de détresse financière. Quelle est la valeur de l'entreprise aujourd'hui, quelle est la valeur de ses fonds propres, quel est son coût des fonds propres et quel est son coût moyen pondéré du capital~? Pourquoi cela ne vaut-il pas environ 7~\%~?
\end{question}

\begin{question}[$\star\star$ \iqroles{Analyste} \iqfirm{hedge fund, banque} \iqlevel{junior, confirmé}]\label{iq:iv2:valuation-and-financial-statements:10}
Le flux de trésorerie disponible est de 100 l'an prochain et croît de 10~\% par an pendant les quatre années suivantes. Après l'année 5, il croît de 3~\% par an à perpétuité. Le coût du capital est de 9~\%, et les flux de chaque année arrivent en fin d'année. Valorisez l'entreprise, et dites quelle part de la valeur provient de la valeur terminale. Que devient la valeur si la croissance terminale est de 4~\%~? Comment la valeur change-t-elle si les flux arrivent au milieu de chaque année~?
\end{question}

\begin{question}[$\star\star$ \iqroles{Analyste} \iqfirm{hedge fund, banque} \iqlevel{junior, confirmé}]\label{iq:iv2:valuation-and-financial-statements:11}
Dans un modèle de valorisation, l'EBITDA de l'année 5 est de 100, le flux de trésorerie disponible de l'année 5 est de 55 et le coût du capital est de 9~\%. Un collègue utilise une valeur terminale égale à 10 fois l'EBITDA. Quel taux de croissance après l'année 5 cela implique-t-il~? Quel multiple impliquerait une croissance de 2~\%~? Un collègue propose 18 fois~: quelle croissance cela exige-t-il, et que lui répondez-vous~?
\end{question}

\begin{question}[$\star\star$ \iqroles{Analyste, QR} \iqfirm{hedge fund, banque} \iqlevel{junior, confirmé}]\label{iq:iv2:valuation-and-financial-statements:12}
Une entreprise gagne 100 après impôt sur ses opérations l'an prochain, et son coût du capital est de 10~\%. Elle croîtra de 5~\% par an à perpétuité en réinvestissant à un rendement constant du capital investi. Quelle est sa valeur si ce rendement est de 5~\%, de 10~\% ou de 20~\%~? Quelle est sa valeur sans croissance~? Dérivez la formule que vous utilisez. Quand la croissance crée-t-elle de la valeur~?
\end{question}

\begin{question}[$\star\star$ \iqroles{Analyste} \iqfirm{hedge fund, banque} \iqlevel{junior, confirmé}]\label{iq:iv2:valuation-and-financial-statements:13}
Deux entreprises ont les mêmes opérations~: un résultat opérationnel de 100, des dotations aux amortissements de 20, un taux d'impôt de 25~\% et une valeur d'entreprise de 1\,200. L'entreprise $A$ n'a pas de dette. L'entreprise $B$ a une dette de 600 à 5~\%. Quels sont, pour chacune, le multiple VE/EBITDA et le multiple cours sur bénéfice~? Pourquoi le second diffère-t-il~? Que vaudrait-il pour $B$ si la dette coûtait 12~\%~?
\end{question}

\begin{question}[$\star\star$ \iqroles{Analyste} \iqfirm{hedge fund, banque} \iqlevel{junior, confirmé}]\label{iq:iv2:valuation-and-financial-statements:14}
Une entreprise sans impôt a un chiffre d'affaires diminué des autres charges de 200 et dépense 100 pour développer un produit. Selon une première politique, elle passe les 100 en charges immédiatement. Selon l'autre, elle les inscrit à l'actif et les amortit linéairement sur 5 ans, à partir de cette année. Pour cette première année, quels sont, dans chaque cas, l'EBITDA, le résultat opérationnel, le flux de trésorerie opérationnel, le flux de trésorerie d'investissement et le flux de trésorerie disponible~? Que devient l'écart de résultat opérationnel si elle dépense 100 chaque année pendant de nombreuses années, et si ses dépenses croissent de 20~\% par an~?
\end{question}

\begin{question}[$\star\star$ \iqroles{Analyste} \iqfirm{hedge fund, banque} \iqlevel{junior}]\label{iq:iv2:valuation-and-financial-statements:15}
Une entreprise réalise un chiffre d'affaires de 1\,000 et un coût des ventes de 600. Elle encaisse auprès de ses clients après 45 jours, conserve son stock pendant 60 jours et paie ses fournisseurs après 30 jours. Prenez 365 jours par an, rapportez les créances au chiffre d'affaires, et les stocks et les dettes fournisseurs au coût des ventes. Quels sont son cycle de conversion de trésorerie et son besoin en fonds de roulement~? Quel montant de trésorerie une croissance de 20~\% absorbe-t-elle si les délais restent les mêmes~? Quel montant de trésorerie une réduction de 15 jours du stock libère-t-elle à la taille actuelle~?
\end{question}

\begin{question}[$\star\star$ \iqroles{Analyste} \iqfirm{hedge fund, banque} \iqlevel{junior, confirmé}]\label{iq:iv2:valuation-and-financial-statements:16}
Une entreprise publie un flux de trésorerie disponible de 150, auquel on réintègre 50 de rémunération en actions, traitée comme un élément sans décaissement. Elle compte 100 actions à 30. Quel rendement du flux de trésorerie disponible la direction affiche-t-elle, et quel serait le rendement si vous traitiez la rémunération en actions comme une charge~? Combien d'actions faudrait-il racheter chaque année pour stabiliser le nombre d'actions~? À 20 fois le flux de trésorerie disponible, quel écart séparerait les deux valorisations~?
\end{question}

\begin{question}[$\star\star$ \iqroles{Analyste} \iqfirm{hedge fund, banque} \iqlevel{junior, confirmé}]\label{iq:iv2:valuation-and-financial-statements:17}
L'entreprise $A$ gagne 100 après impôt, avec 100 actions à 20. Elle rachète l'entreprise $B$, qui gagne 20 après impôt, pour 300. Il n'existe aucune synergie, et le taux d'impôt est de 25~\%. Quel est le nouveau bénéfice par action si elle paie en actions nouvelles~? Quel serait-il si elle emprunte les 300 à 6~\%~? Quelle règle tranche chacun de ces cas~? Que se passe-t-il pour le nombre d'actions si le prix est de 450~? Si une opération relève le bénéfice par action, cela prouve-t-il qu'elle crée de la valeur~?
\end{question}

\begin{question}[$\star\star$ \iqroles{Analyste} \iqfirm{hedge fund, banque} \iqlevel{junior, confirmé}]\label{iq:iv2:valuation-and-financial-statements:18}
Un fonds de LBO paie 10 fois un EBITDA de 100, financé par une dette de 600 et des fonds propres de 400. L'EBITDA croît de 5~\% par an pendant 5 ans, et le flux de trésorerie disponible rembourse 300 de dette sur cette période. Il cède l'entreprise à 10 fois l'EBITDA après 5 ans. Quelles sont la valeur des fonds propres à la sortie, le multiple du capital investi et le rendement annuel~? D'où vient le gain~? Quel est le rendement s'il cède à 9 fois l'EBITDA~?
\end{question}

\begin{question}[$\star\star\star$ \iqroles{Analyste} \iqfirm{hedge fund, banque} \iqlevel{junior, confirmé}]\label{iq:iv2:valuation-and-financial-statements:19}
Une entreprise dont la valeur d'entreprise est de 1\,000 et l'EBITDA de 150 loue un bâtiment pour 100 par an, payés en fin d'année, dans le cadre d'un bail dont il reste 5 ans à courir. Elle traitait jusqu'ici le loyer comme une charge d'exploitation. Elle publie désormais selon une norme qui inscrit le bail au bilan à la valeur actuelle des paiements, actualisés à 6~\%. Quels sont la dette locative et l'actif correspondant~? Quels sont, pour la première année, la dotation aux amortissements, la charge d'intérêts, la variation de l'EBITDA et celle du résultat avant impôt~? Quel est le multiple VE/EBITDA après le changement~? Que devient la trésorerie, et que devient la charge totale sur les 5 ans~?
\end{question}

\begin{question}[$\star\star\star$ \iqroles{Analyste} \iqfirm{hedge fund, banque} \iqlevel{junior, confirmé}]\label{iq:iv2:valuation-and-financial-statements:20}
Une entreprise achète une machine pour 100 et génère 60 par an avant dotations aux amortissements pendant 5 ans. Elle amortit la machine linéairement sur 5 ans dans ses comptes, et linéairement sur 2 ans pour l'impôt. Le taux d'impôt est de 25~\%. Pour chaque année, donnez le résultat net, l'impôt décaissé, le solde d'impôt différé et le flux de trésorerie opérationnel. Quelle est la valeur de ce différé d'impôt à un taux d'actualisation de 8~\%~?
\end{question}

\begin{question}[$\star\star\star$ \iqroles{Analyste} \iqfirm{hedge fund, banque} \iqlevel{junior, confirmé}]\label{iq:iv2:valuation-and-financial-statements:21}
Une entreprise paie 500 en numéraire pour la totalité d'une cible dont les capitaux propres comptables s'élèvent à 200. Les immobilisations corporelles de la cible valent 100 de plus que leur valeur comptable, et ni la réévaluation ni le goodwill ne sont déductibles pour l'impôt. Le taux d'impôt est de 25~\%. Quel goodwill l'acquéreur comptabilise-t-il~? Si la réévaluation est amortie sur 10 ans, comment cela affecte-t-il chaque année le résultat net de l'acquéreur~? Plus tard, le goodwill subit une dépréciation de 100~: que deviennent le résultat net, l'impôt et la trésorerie~?
\end{question}

\begin{question}[$\star\star\star$ \iqroles{Analyste, QR} \iqfirm{hedge fund, banque} \iqlevel{junior, confirmé}]\label{iq:iv2:valuation-and-financial-statements:22}
Une entreprise a une valeur d'entreprise de 3\,000 et vient de dégager un flux de trésorerie disponible de 60. Son coût du capital est de 8~\%. Quel taux de croissance constant à perpétuité le prix implique-t-il~? Si la croissance est au contraire de 3~\% à perpétuité après 10 ans, quel taux de croissance sur ces 10 années le prix implique-t-il~? Quelle est la valeur de l'entreprise si elle croît de 3~\% à partir d'aujourd'hui~? Comment utiliseriez-vous ces chiffres~?
\end{question}

\begin{question}[$\star\star\star$ \iqroles{Analyste} \iqfirm{hedge fund, banque} \iqlevel{junior, confirmé}]\label{iq:iv2:valuation-and-financial-statements:23}
Une enseigne de distribution possède 100 magasins, chacun avec un chiffre d'affaires de 2 et une marge au niveau du magasin de 15~\%, ainsi que des frais centraux de 12. L'an prochain, les ventes des magasins existants augmentent de 3~\%, et 30~\% de toute variation de leur chiffre d'affaires se répercute sur le résultat du magasin, car une grande partie de leurs coûts est fixe. Les frais centraux restent inchangés. Elle ouvre aussi 10 nouveaux magasins au milieu de l'année. Un nouveau magasin réalise un chiffre d'affaires annuel égal à 70~\% de celui d'un magasin mature, avec une marge de 5~\%. Construisez à la main le chiffre d'affaires et le résultat opérationnel de l'an prochain. Quelle est la croissance du résultat opérationnel~? Quel est l'effet d'un point de ventes à magasins comparables sur le résultat, et qu'en est-il si ces ventes sont stables~?
\end{question}

\begin{question}[$\star\star\star\star$ \iqroles{Analyste} \iqfirm{hedge fund} \iqlevel{confirmé}]\label{iq:iv2:valuation-and-financial-statements:24}
Une entreprise publie ces chiffres sur deux ans, et la direction qualifie l'exercice de «~un record~». Le taux d'impôt est de 25~\%. Les actifs s'élevaient à 1\,000 au début de l'année 2 et à 1\,250 à la fin.
\begin{center}\small
\begin{tabular}{@{}lrr@{}}
 & Année 1 & Année 2 \\ \hline
Chiffre d'affaires & 1\,000 & 1\,200 \\
Coût des ventes & 600 & 720 \\
Résultat net & 100 & 130 \\
Flux de trésorerie opérationnel & 110 & 60 \\
Dépenses d'investissement & 40 & 40 \\
Frais de développement immobilisés & 20 & 70 \\
Créances clients & 120 & 210 \\
Stocks & 100 & 170 \\
Produits constatés d'avance & 80 & 60
\end{tabular}
\end{center}
Quelles lignes vous inquiètent~? Calculez ce dont vous avez besoin pour étayer ce constat. De combien le résultat net aurait-il progressé si les frais de développement étaient restés à la part du chiffre d'affaires de l'année 1~? Quelles explications anodines testeriez-vous, et que demanderiez-vous à la direction~?
\end{question}

\begin{question}[$\star\star\star\star$ \iqroles{Analyste, Trader} \iqfirm{hedge fund} \iqlevel{confirmé}]\label{iq:iv2:valuation-and-financial-statements:25}
Vous détenez un titre coté à 90, dimensionné de sorte qu'une baisse jusqu'à votre scénario baissier fasse perdre 1~\% de votre capital. Vous l'avez valorisé à 150 si les marges se redressent, et à 60 sinon, et vous estimez la probabilité d'un redressement à 50~\%. La société abaisse ensuite ses prévisions, et le titre tombe à 63. Vous valorisez désormais les deux scénarios à 120 et à 50. Quelle probabilité de redressement l'ancien cours impliquait-il, et que suppose le nouveau cours~? Quelle probabilité vous faut-il pour obtenir un rendement espéré de 15~\% à partir d'aujourd'hui~? Si vous croyez désormais à 35~\%, quel est votre rendement espéré~? Quelle serait la taille d'une nouvelle position selon la même règle de 1~\%, et combien avez-vous perdu jusqu'ici~? Renforcez-vous, conservez-vous ou réduisez-vous, et qu'est-ce qui vous ferait changer d'avis~?
\end{question}
